\subsection{Limitations}
\label{sec:limits}
We list the limitations that we consider material, from most to least consequential for the claims of this paper.

\paragraph{Compute-bound design} All experiments run on four CPU cores. The image embedding is a frozen ImageNet backbone, the segmentation network is trained at $512\times512$ on 383 images, and we could not train end-to-end grading networks at native resolution, large transformer backbones, retinal foundation models, or heavy detectors. Absolute grading accuracy is therefore lower than what a GPU-trained, fine-tuned system can reach, and the comparison with a fine-tuned ResNet-18 (Section~\ref{sec:res-ft}) is the only end-to-end reference we provide. The \emph{relative} conclusions (evidence tokens versus none, with versus without rule loss) are drawn under identical embeddings, but we cannot exclude that they change when the backbone is fine-tuned or replaced by a stronger one. This is the most important open question for follow-up work.

\paragraph{Single segmentation run} Segmentation and detection results come from one training run per configuration. Differences of a few thousandths between segmentation variants are within the range that seeds could change; we therefore base claims only on the larger differences, and we explicitly avoid significance statements for segmentation.

\paragraph{Train--test discrepancy in lesion evidence} The lesion maps of the 375 grading-training and validation images that are also lesion-training images (Appendix~\ref{app:leak}) are produced by a network that saw their pixel labels, and are therefore cleaner than those of unseen images. This slightly biases the evidence statistics seen during head training towards better lesion maps than at test time (about 4\% of the grading-training images). We removed the corresponding leak from the \emph{test} set, but did not re-train the segmenter on disjoint folds, which would be the cleaner solution.

\paragraph{Rule templates are surrogates} Neovascularisation, intraretinal microvascular abnormalities and venous beading are not annotated in DDR, so rules for severe and proliferative DR are only partial; the zones are image quadrants, not ETDRS fields defined with respect to the optic disc and fovea, because we did not localise those landmarks; and the thresholds are fitted on grade labels, which themselves contain noise. The rules should be read as clinically inspired regularisers and consistency checks, not as an implementation of the grading guideline, and a low violation rate does not by itself prove clinical correctness.

\paragraph{Conformal guarantee} The coverage guarantee is marginal, assumes exchangeability between calibration and test data, and does not hold under distribution shift or conditionally for rare grades (Section~\ref{sec:res-conf}). Because severe DR is rare, class-conditional coverage for grade~3 is estimated from few samples.

\paragraph{External validation} The only external test is a $224\times224$, Gaussian-filtered grading set without lesion annotations. It evaluates robustness of the grading pipeline to a shift in population, camera and resolution, but cannot validate lesion segmentation or detection, and it is not a substitute for prospective multi-centre validation. Conformal calibration was not re-done on the external set.

\paragraph{Labels and annotation} DDR grading labels carry inter-grader noise that we model only through soft ordinal targets. Training-set detection boxes contain duplicates (removed), the detection ground truth is derived from the same lesions as the masks and is therefore not an independent annotation, and microaneurysm annotation is known to be inconsistent. Resizing masks to $512\times512$ with a small area threshold preserves tiny lesions but also dilates some of them by a fraction of a pixel.

\paragraph{No clinical claim} The study is retrospective, on public data. It does not establish clinical utility, safety, or fairness across subgroups; DDR provides no demographic information, so subgroup performance could not be examined.

\subsection{Ethical considerations and responsible use}
Automated DR screening can extend access to care but can also cause harm when deployed without safeguards: false negatives delay treatment of sight-threatening disease, and systematic performance differences across populations, camera types and image quality can widen disparities. Three features of the proposed framework are meant to support responsible use, although none of them is sufficient on its own: lesion-grounded evidence tokens make the basis for a grade inspectable; the rule-violation rate flags internally inconsistent outputs; and conformal sets quantify when the model should abstain from a single-grade decision and refer the case to a human reader. The models and code are released for research only, in line with the non-commercial licence of DDR; any clinical use would require prospective validation, regulatory review, local calibration, and monitoring for distribution shift.

\subsection{Future work}
Four directions follow directly from the limitations. (i) \emph{End-to-end training}: because the evidence bottleneck is differentiable, the segmenter can be fine-tuned from grade supervision, and the fine-tuned image backbone can be replaced by a retinal foundation model \citep{zhou2023retfound}; this requires GPUs. (ii) \emph{Weak lesion supervision}: only 5.5\% of DDR images (757 of 13{,}673) have lesion annotations, but almost all have grades. A teacher--student scheme \citep{tarvainen2017meanteacher} with uncertainty gating \citep{gal2016dropout} and an attention-based multiple-instance loss \citep{ilse2018attentionmil} could use grades as weak lesion supervision; we implemented the building blocks (probability maps, evidence, rules) but did not run this experiment. (iii) \emph{Anatomical zones}: replacing quadrants by ETDRS fields registered to optic disc and fovea would make the rules closer to the clinical definition. (iv) \emph{Prospective and multi-centre evaluation} with conformal re-calibration under shift.
